\documentclass[a4paper,11pt]{book}
%\documentclass[a4paper,twoside,11pt,titlepage]{book}
\usepackage{listings}
\usepackage[utf8]{inputenc}
\usepackage[spanish]{babel}

% \usepackage[style=list, number=none]{glossary} %
%\usepackage{titlesec}
%\usepackage{pailatino}

\decimalpoint
\usepackage{dcolumn}
\newcolumntype{.}{D{.}{\esperiod}{-1}}
\makeatletter
\addto\shorthandsspanish{\let\esperiod\es@period@code}
\makeatother


%\usepackage[chapter]{algorithm}
\RequirePackage{verbatim}
%\RequirePackage[Glenn]{fncychap}
\usepackage{fancyhdr}
\usepackage{graphicx}
\usepackage{afterpage}

\usepackage{longtable}

\usepackage[pdfborder={000}]{hyperref} %referencia

% ********************************************************************
% Re-usable information
% ********************************************************************
\newcommand{\myTitle}{Título del proyecto\xspace}
\newcommand{\myDegree}{Grado en ...\xspace}
\newcommand{\myName}{Nombre Apllido1 Apellido2 (alumno)\xspace}
\newcommand{\myProf}{Nombre Apllido1 Apellido2 (tutor1)\xspace}
\newcommand{\myOtherProf}{Nombre Apllido1 Apellido2 (tutor2)\xspace}
%\newcommand{\mySupervisor}{Put name here\xspace}
\newcommand{\myFaculty}{Escuela Técnica Superior de Ingenierías Informática y de
Telecomunicación\xspace}
\newcommand{\myFacultyShort}{E.T.S. de Ingenierías Informática y de
Telecomunicación\xspace}
\newcommand{\myDepartment}{Departamento de ...\xspace}
\newcommand{\myUni}{\protect{Universidad de Granada}\xspace}
\newcommand{\myLocation}{Granada\xspace}
\newcommand{\myTime}{\today\xspace}
\newcommand{\myVersion}{Version 0.1\xspace}


\hypersetup{
pdfauthor = {\myName (email (en) ugr (punto) es)},
pdftitle = {\myTitle},
pdfsubject = {},
pdfkeywords = {palabra_clave1, palabra_clave2, palabra_clave3, ...},
pdfcreator = {LaTeX con el paquete ....},
pdfproducer = {pdflatex}
}

%\hyphenation{}


%\usepackage{doxygen/doxygen}
%\usepackage{pdfpages}
\usepackage{url}
\usepackage{colortbl,longtable}
\usepackage[stable]{footmisc}
%\usepackage{index}

%\makeindex
%\usepackage[style=long, cols=2,border=plain,toc=true,number=none]{glossary}
% \makeglossary

% Definición de comandos que me son tiles:
%\renewcommand{\indexname}{Índice alfabético}
%\renewcommand{\glossaryname}{Glosario}

\pagestyle{fancy}
\fancyhf{}
\fancyhead[LO]{\leftmark}
\fancyhead[RE]{\rightmark}
\fancyhead[RO,LE]{\textbf{\thepage}}
\renewcommand{\chaptermark}[1]{\markboth{\textbf{#1}}{}}
\renewcommand{\sectionmark}[1]{\markright{\textbf{\thesection. #1}}}

\setlength{\headheight}{1.5\headheight}

\newcommand{\HRule}{\rule{\linewidth}{0.5mm}}
%Definimos los tipos teorema, ejemplo y definición podremos usar estos tipos
%simplemente poniendo \begin{teorema} \end{teorema} ...
\newtheorem{teorema}{Teorema}[chapter]
\newtheorem{ejemplo}{Ejemplo}[chapter]
\newtheorem{definicion}{Definición}[chapter]

\definecolor{gray97}{gray}{.97}
\definecolor{gray75}{gray}{.75}
\definecolor{gray45}{gray}{.45}
\definecolor{gray30}{gray}{.94}

\lstset{ frame=Ltb,
     framerule=0.5pt,
     aboveskip=0.5cm,
     framextopmargin=3pt,
     framexbottommargin=3pt,
     framexleftmargin=0.1cm,
     framesep=0pt,
     rulesep=.4pt,
     backgroundcolor=\color{gray97},
     rulesepcolor=\color{black},
     %
     stringstyle=\ttfamily,
     showstringspaces = false,
     basicstyle=\scriptsize\ttfamily,
     commentstyle=\color{gray45},
     keywordstyle=\bfseries,
     %
     numbers=left,
     numbersep=6pt,
     numberstyle=\tiny,
     numberfirstline = false,
     breaklines=true,
   }
 
% minimizar fragmentado de listados
\lstnewenvironment{listing}[1][]
   {\lstset{#1}\pagebreak[0]}{\pagebreak[0]}

\lstdefinestyle{CodigoC}
   {
	basicstyle=\scriptsize,
	frame=single,
	language=C,
	numbers=left
   }
\lstdefinestyle{CodigoC++}
   {
	basicstyle=\small,
	frame=single,
	backgroundcolor=\color{gray30},
	language=C++,
	numbers=left
   }

 
\lstdefinestyle{Consola}
   {basicstyle=\scriptsize\bf\ttfamily,
    backgroundcolor=\color{gray30},
    frame=single,
    numbers=none
   }


\newcommand{\bigrule}{\titlerule[0.5mm]}


%Para conseguir que en las páginas en blanco no ponga cabecerass
\makeatletter
\def\clearpage{%
  \ifvmode
    \ifnum \@dbltopnum =\m@ne
      \ifdim \pagetotal <\topskip
        \hbox{}
      \fi
    \fi
  \fi
  \newpage
  \thispagestyle{empty}
  \write\m@ne{}
  \vbox{}
  \penalty -\@Mi
}
\makeatother

\usepackage{pdfpages}
\begin{document}
\begin{titlepage}
 
 
    \newlength{\centeroffset}
    \setlength{\centeroffset}{-0.5\oddsidemargin}
    \addtolength{\centeroffset}{0.5\evensidemargin}
    \thispagestyle{empty}
    
    \noindent\hspace*{\centeroffset}\begin{minipage}{\textwidth}
    
    \centering
    \includegraphics[width=0.9\textwidth]{imagenes/logo_ugr.jpg}\\[1.4cm]
    
    \textsc{ \Large TRABAJO FIN DE GRADO\\[0.2cm]}
    \textsc{ INGENIERÍA INFORMÁTICA}\\[1cm]
    % Upper part of the page
    % 
    % Title
    {\Huge\bfseries Titulo del Proyecto\\
    }
    \noindent\rule[-1ex]{\textwidth}{3pt}\\[3.5ex]
    {\large\bfseries Subtitulo del Proyecto}
    \end{minipage}
    
    \vspace{2.5cm}
    \noindent\hspace*{\centeroffset}\begin{minipage}{\textwidth}
    \centering
    
    \textbf{Autor}\\ {Ángel Sánchez Guerrero}\\[2.5ex]
    \textbf{Directores}\\
    {Nombre Apellido1 Apellido2 (tutor1)\\
    Nombre Apellido1 Apellido2 (tutor2)}\\[2cm]
    \includegraphics[width=0.3\textwidth]{imagenes/etsiit_logo.png}\\[0.1cm]
    \textsc{Escuela Técnica Superior de Ingenierías Informática y de Telecomunicación}\\
    \textsc{---}\\
    Granada, Junio de 2025
    \end{minipage}
    %\addtolength{\textwidth}{\centeroffset}
    %\vspace{\stretch{2}}
    \end{titlepage}
    
    
    
\chapter*{}
%\thispagestyle{empty}
%\cleardoublepage

%\thispagestyle{empty}

\begin{titlepage}
 
 
    \setlength{\centeroffset}{-0.5\oddsidemargin}
    \addtolength{\centeroffset}{0.5\evensidemargin}
    \thispagestyle{empty}
    
    \noindent\hspace*{\centeroffset}\begin{minipage}{\textwidth}
    
    \centering
    %\includegraphics[width=0.9\textwidth]{imagenes/logo_ugr.jpg}\\[1.4cm]
    
    %\textsc{ \Large PROYECTO FIN DE CARRERA\\[0.2cm]}
    %\textsc{ INGENIERÍA EN INFORMÁTICA}\\[1cm]
    % Upper part of the page
    % 
    
     \vspace{3.3cm}
    
    %si el proyecto tiene logo poner aquí
    %\includegraphics{imagenes/logo.png} 
    % \vspace{0.5cm}
    
    % Title
    
    {\Huge\bfseries Título del proyecto\\
    }
    \noindent\rule[-1ex]{\textwidth}{3pt}\\[3.5ex]
    {\large\bfseries Subtítulo del proyecto.\\[4cm]}
    \end{minipage}
    
    \vspace{2.5cm}
    \noindent\hspace*{\centeroffset}\begin{minipage}{\textwidth}
    \centering
    
    \textbf{Autor}\\ {Ángel Sánchez Guerrero}\\[2.5ex]
    \textbf{Directores}\\
    {Nombre Apellido1 Apellido2 (tutor1)\\
    Nombre Apellido1 Apellido2 (tutor2)}\\[2cm]
    %\includegraphics[width=0.15\textwidth]{imagenes/tstc.png}\\[0.1cm]
    %\textsc{Departamento de Teoría de la Señal, Telemática y Comunicaciones}\\
    %\textsc{---}\\
    %Granada, mes de 201
    \end{minipage}
    %\addtolength{\textwidth}{\centeroffset}
    \vspace{\stretch{2}}
    
     
    \end{titlepage}
    
    
    



\cleardoublepage
\thispagestyle{empty}

\begin{center}
{\large\bfseries Título del Proyecto: Subtítulo del proyecto}\\
\end{center}
\begin{center}
Ángel Sánchez Guerrero\\
\end{center}

%\vspace{0.7cm}
\noindent{\textbf{Palabras clave}: palabra\_clave1, palabra\_clave2, palabra\_clave3, ......}\\

\vspace{0.7cm}
\noindent{\textbf{Resumen}}\\

Poner aquí el resumen.
\cleardoublepage


\thispagestyle{empty}

\begin{center}
{\large\bfseries Project Title: Project Subtitle}\\
\end{center}
\begin{center}
Ángel Sánchez Guerrero\\
\end{center}

%\vspace{0.7cm}
\noindent{\textbf{Keywords}: Keyword1, Keyword2, Keyword3, ....}\\

\vspace{0.7cm}
\noindent{\textbf{Abstract}}\\

Write here the abstract in English.

\chapter*{}
\thispagestyle{empty}

\noindent\rule[-1ex]{\textwidth}{2pt}\\[4.5ex]

Yo, \textbf{Ángel Sánchez Guerrero}, alumno de la titulación INGENIERÍA INFORMÁTICA de la \textbf{Escuela Técnica Superior
de Ingenierías Informática y de Telecomunicación de la Universidad de Granada}, con DNI 77147702K, autorizo la
ubicación de la siguiente copia de mi Trabajo Fin de Grado en la biblioteca del centro para que pueda ser
consultada por las personas que lo deseen.

\vspace{6cm}

\noindent Fdo: Ángel Sánchez Guerrero

\vspace{2cm}

\begin{flushright}
Granada a X de mes de 201 .
\end{flushright}


\chapter*{}
\thispagestyle{empty}

\noindent\rule[-1ex]{\textwidth}{2pt}\\[4.5ex]

D. \textbf{Nombre Apellido1 Apellido2 (tutor1)}, Profesor del Área de XXXX del Departamento YYYY de la Universidad de Granada.

\vspace{0.5cm}

D. \textbf{Nombre Apellido1 Apellido2 (tutor2)}, Profesor del Área de XXXX del Departamento YYYY de la Universidad de Granada.


\vspace{0.5cm}

\textbf{Informan:}

\vspace{0.5cm}

Que el presente trabajo, titulado \textit{\textbf{Título del proyecto, Subtítulo del proyecto}},
ha sido realizado bajo su supervisión por \textbf{Nombre Apellido1 Apellido2 (alumno)}, y autorizamos la defensa de dicho trabajo ante el tribunal
que corresponda.

\vspace{0.5cm}

Y para que conste, expiden y firman el presente informe en Granada a X de mes de 201 .

\vspace{1cm}

\textbf{Los directores:}

\vspace{5cm}

\noindent \textbf{Nombre Apellido1 Apellido2 (tutor1) \ \ \ \ \ Nombre Apellido1 Apellido2 (tutor2)}

\chapter*{Agradecimientos}
\thispagestyle{empty}

       \vspace{1cm}


Poner aquí agradecimientos...


\frontmatter
\tableofcontents
\listoffigures
\listoftables
%
\mainmatter
\setlength{\parskip}{5pt}

\chapter{Introducción}

En este primer capítulo se expone el contexto y motivación del proyecto, los objetivos planteados, y la estructura que se seguirá en la memoria.

\section{Contexto y motivación}

La transformación digital del sector sanitario ha avanzado de forma decidida en las últimas décadas. En Estados Unidos, por ejemplo, cerca del 96\% de los hospitales ya habían adoptado sistemas de historiales clínicos electrónicos (EHR) para 2015 \cite{Henry2016_EHR}. Este avance ha generado un volumen masivo de datos clínicos, fundamentales para la investigación en áreas como la epidemiología o el desarrollo de modelos predictivos \cite{Halevy2009_data}. Sin embargo, aunque la generación de datos es continua y creciente, la mayoría de estos registros están sujetos a estrictas restricciones de acceso debido a cuestiones legales, éticas y, sobre todo, de privacidad del paciente. Esta protección, imprescindible para salvaguardar la confidencialidad, limita el potencial de los datos para la investigación y la innovación clínica.

En este contexto, la aparición de bases de datos como MIMIC-III \cite{MIMICIII_paper} (Medical Information Mart for Intensive Care) supuso un hito al demostrar que es posible compartir datos clínicos reales de forma segura mediante procesos rigurosos de desidentificación, permitiendo así la realización de cientos de estudios que han mejorado la práctica clínica \cite{Kallout2025_contribution}. El proyecto MIMIC-IV \cite{MIMICIV_paper, MIMICIV_dataset}, sobre el que se fundamenta este trabajo, continúa y amplía ese legado con información clínica detallada y anonimizada del Beth Israel Deaconess Medical Center de Boston. En su versión 3.1, MIMIC-IV incluye datos de más de 364.000 pacientes y casi 550.000 admisiones hospitalarias, abarcando desde datos demográficos y resultados de laboratorio hasta procedimientos médicos y evolución clínica.

\newpage
No obstante, a pesar de la riqueza y el potencial de MIMIC-IV, persiste un reto relevante: extraer conocimiento útil de estos datos sigue siendo complejo para quienes no poseen formación técnica avanzada. La barrera tecnológica impide que muchos profesionales sanitarios, investigadores clínicos y estudiantes puedan aprovechar plenamente la información disponible. Ante esta situación, la motivación fundamental de este proyecto es precisamente derribar ese obstáculo, desarrollando una plataforma donde cualquier usuario pueda visualizar, explorar y descubrir patrones complejos en los datos, acelerando así la generación de hipótesis y la toma de decisiones fundamentadas en evidencia real.

\section{Objetivos}

\subsection{Objetivo general}

Este proyecto tiene como objetivo principal desarrollar una aplicación web con Inteligencia Artificial que permita visualizar, analizar, predecir y explicar historiales clínicos contenidos en MIMIC-IV.

\subsection{Objetivos específicos}

\begin{itemize}
\item \textbf{OE1.} Almacenar el conjunto de datos MIMIC-IV en una base de datos para poder realizar consultas eficientes.
\item \textbf{OE2.} Desarrollar un backend API RESTful que sirva como intermediario entre la base de datos y el frontend.
\item \textbf{OE3.} Construir una plataforma web con interfaz intuitiva y visualizaciones dinámicas para facilitar la interpretación de los datos.
\item \textbf{OE4.} Implementar un asistente de IA que permita realizar consultas en lenguaje natural sobre la base de datos.
\item \textbf{OE5.} Implementar la predicción de datos con IA para obtener indicadores de riesgo.
\item \textbf{OE6.} Configurar el entorno de producción y evaluar la usabilidad, el rendimiento y la utilidad de la plataforma final.
\end{itemize}

\newpage
\section{Estructura de la memoria}

El presente documento se estructura de la siguiente manera:

\begin{itemize}
    \item \textbf{Capítulo 1: Introducción}. Presenta el contexto, motivación, objetivos y la estructura del documento.
    
    \item \textbf{Capítulo 2: Estado del arte}. Revisión de la literatura y trabajos previos.

    ...
    
    %\item \textbf{Capítulo 3: Diseño}. Descripción detallada de la arquitectura, modelos de datos y diseño de interfaces de la aplicación.
    
    %\item \textbf{Capítulo 4: Implementación}. Explicación del proceso de desarrollo de la aplicación, incluyendo tecnologías utilizadas, configuración del entorno y aspectos técnicos relevantes.
    
    %\item \textbf{Capítulo 5: Resultados y evaluación}. Presentación de los resultados obtenidos y evaluación del desempeño de la aplicación en términos de usabilidad, eficiencia y utilidad.
    
    %\item \textbf{Capítulo 6: Conclusiones y trabajo futuro}. Resumen de las contribuciones del proyecto, limitaciones encontradas y posibles líneas de trabajo futuro.
\end{itemize}


\chapter{Estado del arte}

En este capítulo se analizan los fundamentos teóricos y las tecnologías existentes que conforman la base de este proyecto. Se revisan las principales bases de datos clínicas abiertas, las arquitecturas de software para la gestión de datos sanitarios, las herramientas de visualización interactiva y las aplicaciones de la Inteligencia Artificial en el ámbito clínico.

\section{El Movimiento hacia la Ciencia Abierta en la Investigación Médica}

En las últimas décadas se ha consolidado un movimiento hacia los datos clínicos abiertos, promoviendo la disponibilidad de bases de datos sanitarias para la comunidad investigadora. La idea de ciencia abierta sostiene que la libre disponibilidad de datos y códigos favorece la transparencia, la reproducibilidad y la colaboración en la investigación biomédica \cite{Lvovs2025_balancing}. Sin embargo, históricamente los datos clínicos han estado protegidos en archivos hospitalarios, difíciles de acceder por razones legales, éticas y técnicas. En respuesta, iniciativas académicas y gubernamentales han impulsado la creación de conjuntos de datos clínicos de acceso abierto, debidamente anonimizados, que permiten a investigadores de todo el mundo analizar información real de pacientes sin vulnerar la privacidad. Estos recursos han transformado la investigación médica al reducir barreras de acceso y fomentar prácticas reproducibles en análisis de datos sanitarios

\newpage
Los datos clínicos se pueden clasificar en cuatro grandes modalidades complementarias:


(i) \emph{Señales fisiológicas de alta frecuencia.} Series temporales de ECG, presión arterial invasiva, fotopletismografía o EEG. Podemos destacar la clásica MIT-BIH Arrhythmia Database (1980), patrón de referencia para algoritmos de ECG \cite{Impact_MIT-BIH}; el MIMIC Waveform Database, que enlaza miles de horas de señales con los datos clínicos de MIMIC \cite{Moody2022_MIMICIVWaveform}; y HiRID, con 712 variables registradas cada dos minutos en más de 30 000 estancias UCI \cite{Faltys2021HiRID}.

(ii) \emph{Imágenes médicas con anotaciones diagnósticas.}  Abarcan radiografías, TC, RM, PET e incluso histopatología digital, cada una acompañada de etiquetas o informes de hallazgo.  Las bases más usadas en radiología son ChestXray14 y CheXpert, ambas con cientos de miles de radiografías torácicas etiquetadas para patologías pulmonares \cite{irvin2019chexpertlargechestradiograph}, y MIMIC-CXR por parte de la familia MIMIC \cite{Johnson2019_MIMICCXR}.  En oncología destacan las colecciones TCGA/TCIA \cite{TCGA,Clark2013_TCIA}, y en neuroimagen la iniciativa ADNI \cite{Petersen2010_ADNI}.  Estos recursos posibilitan evaluar modelos de visión computacional con criterios homogéneos.

(iii) \emph{Texto clínico desidentificado.}  Incluye notas de evolución, informes radiológicos, resúmenes de alta, entre otros.  Representa alrededor del 80 \% de la información clínica, pero su liberación es más compleja por contener PHI (Protected Health Information).  Los desafíos i2b2/n2c2, pioneros al proveer corpora anonimizados, constituyen la plataforma estándar para comparar sistemas de procesamiento del lenguaje natural médico \cite{n2c2}.  El ecosistema MIMIC también complementa este dominio con millones de notas clínicas \cite{Johnson2023_MIMICIVNote}. 

(iv) \emph{Datos estructurados de historias clínicas electrónicas (EHR).}  Comprenden tablas con demografía, diagnósticos (ICD-9/10), procedimientos, resultados de laboratorio, medicación y constantes vitales.  Son la base de estudios epidemiológicos y de la construcción de modelos pronósticos.  Entre los conjuntos abiertos más influyentes figuran NHANES \cite{NHANES}, UK Biobank \cite{Sudlow2015_UKBiobank}, eICU \cite{Pollard2018} y, sobre todo, la serie MIMIC en la que se enfoca este trabajo.

Desde 1999 PhysioNet \cite{PhysioNet_paper} estableció un marco seguro y reproducible para compartir todo tipo de registros hospitalarios.  Bajo ese paraguas nació el ecosistema MIMIC: la primera versión (1996) contenía 90 pacientes UCI \cite{Moody1996_MIMIC}; MIMIC-II (2011) multiplicó tamaño y variables al extraer directamente de los sistemas clínicos \cite{Saeed2011_MIMICII}; MIMIC-III (2015) superó los 40 000 pacientes y se convirtió en la referencia mundial \cite{MIMICIII_paper}; MIMIC-IV \cite{MIMICIV_paper} amplía y moderniza el conjunto.  Esta trayectoria demuestra cómo los EHR han pasado a ser un recurso científico global, permitiendo investigar la fisiopatología crítica con una profundidad antes impensable.

\newpage
En conjunto, el panorama de datos clínicos abiertos ofrece múltiples alternativas, pero MIMIC-IV destaca como la más idónea para este trabajo: combina un gran volumen de datos reciente, buena documentación, y módulos complementarios. Sobre esta base, el proyecto diseñará una plataforma que haga de este conjunto de datos una fuente accesible de conocimiento, alineada con los principios de ciencia abierta que han guiado todo este repaso bibliográfico.


\section{Tecnologías para la Gestión y Almacenamiento de Datos Sanitarios}

En esta sección justificamos la arquitectura tecnológica del backend.
Bases de Datos NoSQL vs. SQL: Analizar por qué las bases de datos NoSQL como MongoDB son especialmente adecuadas para datos clínicos, que son heterogéneos y semi-estructurados, en comparación con las bases de datos relacionales tradicionales.
Arquitecturas de Backend y APIs: Hablar sobre la tendencia de usar arquitecturas orientadas a servicios (como mencionaba tu tutor) y el papel de las APIs RESTful (usando FastAPI como ejemplo) para desacoplar el frontend del backend y servir datos de manera eficiente.

\section{Herramientas y Técnicas de Visualización Interactiva de Datos}

Esta parte se centra en cómo vamos a "ver" los datos.
Paradigmas de Visualización: Explicar la diferencia entre visualizaciones estáticas y dinámicas/interactivas, y por qué estas últimas son cruciales para la exploración de datos complejos.
Librerías y Plataformas: Analizar las herramientas más potentes del mercado. Mencionar D3.js como la librería de bajo nivel que ofrece máxima flexibilidad, y herramientas de más alto nivel como Observable Plot (la que usaremos), Plotly, o incluso plataformas como Tableau, explicando las ventajas de nuestra elección para este proyecto.

\section{Aplicaciones de Inteligencia Artificial en el Análisis Clínico}

Aquí abordamos la parte más innovadora del proyecto. Podemos dividirla en dos sub-apartados que se corresponden con tus objetivos:
Modelos Predictivos para Indicadores de Riesgo: Revisar el estado del arte en la aplicación de modelos de machine learning (desde regresiones hasta redes neuronales) para predecir eventos clínicos como mortalidad, reingresos, etc., usando datos de EHR.
Asistentes Conversacionales y Procesamiento de Lenguaje Natural (NLP): Hablar sobre la tendencia emergente de usar Modelos de Lenguaje Grandes (LLMs) para crear interfaces conversacionales que permitan a usuarios no técnicos "preguntar" a las bases de datos en lenguaje natural, justificando así la funcionalidad del chat con IA.
También RAG vs MCP

%
%\input{capitulos/02_EspecificacionRequisitos}
%
%\input{capitulos/03_Planificacion}
%
%\input{capitulos/04_Analisis}
%
%\input{capitulos/05_Diseno}
%
%\input{capitulos/06_Implementacion}
%
%\input{capitulos/07_Pruebas}
%
%\input{capitulos/08_Conclusiones}
%
%%\chapter{Conclusiones y Trabajos Futuros}
%
%
%\nocite{*}
\bibliography{bibliografia/bibliografia}\addcontentsline{toc}{chapter}{Bibliografía}
\bibliographystyle{unsrturl}
%
%\appendix
%\input{apendices/manual_usuario/manual_usuario}
%%\input{apendices/paper/paper}
%\input{glosario/entradas_glosario}
% \addcontentsline{toc}{chapter}{Glosario}
% \printglossary
\chapter*{}
\thispagestyle{empty}

\end{document}
